\chapter{Amenazas a la validez y limitaciones}
\label{cap:amenazas-limitaciones-final}

Este cap\'itulo integra y actualiza, con la evidencia real de v1.0.0, el
borrador de Jaime (\texttt{docs/informe/borradores/jaime/metodologia-y-amenazas.md}),
sin minimizar ninguna amenaza.

\section{Validez interna}

\begin{itemize}[nosep]
  \item Todas las mediciones de rendimiento y seguridad locales se
    ejecutaron en una sola m\'aquina de desarrollo, sin control de
    interferencia de otros procesos del sistema operativo.
  \item El certificado TLS del entorno de desarrollo es autofirmado; en
    producci\'on, HTTPS lo gestiona Render autom\'aticamente, fuera del
    control directo del equipo.
  \item El hit ratio de cach\'e se puede medir con dos m\'etodos que
    producen cifras distintas (global vs. aislado por clave); se document\'o
    expl\'icitamente cu\'al se us\'o y por qu\'e.
  \item El an\'alisis Wilcoxon/Rosenthal muestra tama\~nos de efecto no
    uniformes entre corridas ($r$ de $-0{,}07$ a $-0{,}84$), lo que indica
    variabilidad de entorno no completamente controlada entre corridas
    consecutivas.
\end{itemize}

\section{Validez externa}

\begin{itemize}[nosep]
  \item Todas las mediciones locales se realizaron en un entorno
    acad\'emico; el \'unico entorno equivalente a producci\'on es el
    despliegue real en el plan gratuito de Render, que tiene sus propias
    limitaciones (secci\'on~\ref{sec:limitacion-30-dias}).
  \item La muestra SUS ($n=18$, reclutada por conveniencia) no es
    necesariamente representativa de usuarios finales de una cl\'inica
    veterinaria real; un tama\~no muestral mayor reduce el margen del
    intervalo de confianza pero no cambia el m\'etodo de reclutamiento.
  \item Lighthouse depende del entorno de ejecuci\'on (hardware, red,
    versi\'on del navegador); sus puntajes no son un valor absoluto
    reproducible en cualquier m\'aquina.
  \item Los usuarios de prueba del backend (cuentas acad\'emicas, admin
    sembrado) no son usuarios reales de una cl\'inica veterinaria.
\end{itemize}

\section{Validez de constructo}

\begin{itemize}[nosep]
  \item JaCoCo mide cobertura de ejecuci\'on, no ausencia de defectos.
  \item Lighthouse no representa toda la experiencia de usuario real; se
    complementa con SUS como medici\'on independiente, pero ninguna
    sustituye una evaluaci\'on de UX cualitativa dedicada.
  \item SUS mide percepci\'on de usabilidad, no eficacia cl\'inica ni
    operacional real en un contexto veterinario.
  \item ZAP y SpotBugs no prueban la ausencia total de vulnerabilidades,
    solo la ausencia de los patrones que sus reglas configuradas detectan.
  \item Las pruebas de integraci\'on con Testcontainers verifican
    comportamiento real de PostgreSQL para las rutinas nativas; la mayor\'ia
    del resto de la suite usa H2 en memoria, cuyo comportamiento no es
    id\'entico al de PostgreSQL real.
\end{itemize}

\section{Validez de conclusi\'on}

\begin{itemize}[nosep]
  \item Cada corrida de k6 ($\sim$3\,200 peticiones) corresponde a una
    ventana de tiempo corta ($\sim$30--35~s); no se midieron tendencias en
    periodos prolongados.
  \item El tama\~no muestral de SUS ($n=18$) sigue siendo insuficiente para
    an\'alisis de subgrupos demogr\'aficos con potencia estad\'istica
    adecuada.
  \item Ninguna medici\'on de este informe implica dise\~no experimental
    con grupos de control; ninguna mejora observada entre versiones se
    atribuye causalmente a un cambio espec\'ifico.
  \item Las 205 pruebas y la cobertura JaCoCo se ejecutan en un ambiente
    de prueba mixto (H2 + Testcontainers/PostgreSQL para las rutinas
    nativas), no en un \'unico motor de base de datos consistente en toda
    la suite.
\end{itemize}

\section{Limitaci\'on temporal expl\'icita: operaci\'on sostenida no verificable hoy}
\label{sec:limitacion-30-dias}

El plan gratuito de PostgreSQL en Render \textbf{elimina los datos a los 30
d\'ias} de creado el recurso, documentado expl\'icitamente en
\texttt{docs/despliegue/DEPLOYMENT.md} y en la pol\'itica de retenci\'on de
respaldos de \texttt{docs/despliegue/BACKUP.md} (retenci\'on de 30 d\'ias,
purgado autom\'atico de respaldos con m\'as de 31 d\'ias). Esto significa
que, a la fecha de cierre de este informe, \textbf{no es posible verificar
emp\'iricamente una condici\'on de operaci\'on sostenida del sistema en
producci\'on durante 30 d\'ias o m\'as}: el despliegue actual es real y
verificable (\emph{healthcheck} \texttt{UP}, HTTPS activo), pero su
horizonte temporal de evidencia continua est\'a acotado por la propia
pol\'itica del plan gratuito usado. Este informe documenta esta condici\'on
como una \textbf{limitaci\'on temporal genuina}, no como un requisito
cumplido ni como uno incumplido por negligencia: verificarla exigir\'ia
migrar a un plan pago de base de datos persistente (o instrumentar un
respaldo/restauraci\'on autom\'atico verificado en producci\'on real durante
ese periodo), decisi\'on que excede el alcance de esta entrega acad\'emica.

\section{Otras limitaciones del estudio}

\begin{itemize}[nosep]
  \item \textbf{Contexto exclusivamente acad\'emico:} ninguna medici\'on
    proviene de un despliegue operativo con usuarios reales de una
    cl\'inica veterinaria.
  \item \textbf{Sin evaluaci\'on industrial ni con desarrolladores
    profesionales externos.}
  \item \textbf{Sin comparaci\'on emp\'irica directa con sistemas de
    gesti\'on veterinaria alternativos:} toda comparaci\'on es contra
    umbrales t\'ecnicos propios o contra la literatura, no contra artefactos
    competidores evaluados en las mismas condiciones.
  \item \textbf{Lighthouse: mecanismo de sesi\'on no documentado en
    prosa.} La corrida real del 2026-08-18 (mobile+desktop, 12 corridas,
    4/4 categor\'ias cumplidas, incluido SEO~100/100) audit\'o
    \texttt{/mascotas} directamente (\texttt{finalUrl} sin redirecci\'on a
    \texttt{/login} en los seis reportes de esa ruta, a diferencia de la
    corrida de agosto); sin embargo, el equipo no document\'o
    expl\'icitamente en \texttt{docs/mediciones/lighthouse/README.md} qu\'e
    mecanismo permiti\'o evitar la redirecci\'on del \emph{guard} de
    autenticaci\'on. Este informe reporta el dato num\'erico real
    (verificado en los JSON crudos) sin afirmar un mecanismo de sesi\'on no
    documentado por el equipo. La narrativa completa del
    \texttt{README.md} de esa carpeta a\'un no fue actualizada para
    describir la corrida del 2026-08-18 en prosa.
  \item \textbf{Rate limiting de login no distribuido:} en memoria, por
    instancia \'unica del backend.
  \item \textbf{Sin integraci\'on con SIEM} para los logs de auditor\'ia.
  \item \textbf{Requisitos heredados de la Entrega~1A a\'un pendientes:}
    historial cl\'inico longitudinal completo, calendario interactivo de
    citas, IoT, facturaci\'on digital, reportes exportables y
    recomendaciones por IA (cap\'itulo~\ref{cap:requisitos}).
  \item \textbf{Bloque Jaime de PRISMA sin conteo de exclusiones:} los 10
    estudios incluidos (meta de $\geq$8--9 cumplida) est\'an documentados y
    el diagrama de flujo ya fue regenerado con los tres bloques
    cuantificados; el \'unico gap que persiste es que Jaime no document\'o
    su conteo de candidatos evaluados/excluidos, se\~nalado con l\'inea
    punteada en el propio diagrama (cap\'itulo~\ref{cap:trabajos-relacionados}).
  \item \textbf{Tag \texttt{v1.0.0} y GitHub Release todav\'ia no creados}
    (secci\'on~\ref{sec:pendientes-honestos-publicacion}), fuera del
    alcance expl\'icito de esta tarea de redacci\'on del informe.
\end{itemize}

Ninguna de estas limitaciones se presenta como resuelta; todas condicionan
expl\'icitamente el alcance de las conclusiones del
cap\'itulo~\ref{cap:conclusiones-final}.
